\documentclass[11pt]{article}
% The preamble every project's paper shares: % The preamble every project's paper shares: % The preamble every project's paper shares: \input{../../preamble} from
% projects/<project>/paper/. Boards come from chessboard and xskak, plots from
% pgfplots; every number a paper prints comes from its generated.tex.
\usepackage[LSBC4,T1]{fontenc}
\usepackage[utf8]{inputenc}
\usepackage{lmodern}
\usepackage[margin=1in]{geometry}
\usepackage{microtype}
\usepackage{amsmath}
\usepackage{booktabs}
\usepackage{tabularx}
\usepackage{array}
\usepackage{xcolor}
\usepackage[obeyspaces,spaces]{url}
\usepackage{xskak}
\usepackage{pgfplots}
\pgfplotsset{compat=1.18}
\usepackage{caption}
\usepackage{subcaption}
\usepackage{enumitem}
\usepackage[section]{placeins}
\usepackage[hidelinks]{hyperref}

\setchessboard{boardfontencoding=LSBC4, boardfontfamily=skaknew, blackfieldcolor=black!22,
  whitefieldcolor=white, setfontcolors, showmover=false, label=true, labelfontsize=6pt,
  boardfontsize=14pt}

% A move or line in UCI, and a position in FEN, both breakable at a slash.
\newcommand{\uci}[1]{\texttt{#1}}
\newcommand{\fen}[1]{{\footnotesize\path{#1}}}
\newcolumntype{L}{>{\raggedright\arraybackslash}X}

\setlength{\parskip}{0.4em}
\setlength{\parindent}{0pt}
 from
% projects/<project>/paper/. Boards come from chessboard and xskak, plots from
% pgfplots; every number a paper prints comes from its generated.tex.
\usepackage[LSBC4,T1]{fontenc}
\usepackage[utf8]{inputenc}
\usepackage{lmodern}
\usepackage[margin=1in]{geometry}
\usepackage{microtype}
\usepackage{amsmath}
\usepackage{booktabs}
\usepackage{tabularx}
\usepackage{array}
\usepackage{xcolor}
\usepackage[obeyspaces,spaces]{url}
\usepackage{xskak}
\usepackage{pgfplots}
\pgfplotsset{compat=1.18}
\usepackage{caption}
\usepackage{subcaption}
\usepackage{enumitem}
\usepackage[section]{placeins}
\usepackage[hidelinks]{hyperref}

\setchessboard{boardfontencoding=LSBC4, boardfontfamily=skaknew, blackfieldcolor=black!22,
  whitefieldcolor=white, setfontcolors, showmover=false, label=true, labelfontsize=6pt,
  boardfontsize=14pt}

% A move or line in UCI, and a position in FEN, both breakable at a slash.
\newcommand{\uci}[1]{\texttt{#1}}
\newcommand{\fen}[1]{{\footnotesize\path{#1}}}
\newcolumntype{L}{>{\raggedright\arraybackslash}X}

\setlength{\parskip}{0.4em}
\setlength{\parindent}{0pt}
 from
% projects/<project>/paper/. Boards come from chessboard and xskak, plots from
% pgfplots; every number a paper prints comes from its generated.tex.
\usepackage[LSBC4,T1]{fontenc}
\usepackage[utf8]{inputenc}
\usepackage{lmodern}
\usepackage[margin=1in]{geometry}
\usepackage{microtype}
\usepackage{amsmath}
\usepackage{booktabs}
\usepackage{tabularx}
\usepackage{array}
\usepackage{xcolor}
\usepackage[obeyspaces,spaces]{url}
\usepackage{xskak}
\usepackage{pgfplots}
\pgfplotsset{compat=1.18}
\usepackage{caption}
\usepackage{subcaption}
\usepackage{enumitem}
\usepackage[section]{placeins}
\usepackage[hidelinks]{hyperref}

\setchessboard{boardfontencoding=LSBC4, boardfontfamily=skaknew, blackfieldcolor=black!22,
  whitefieldcolor=white, setfontcolors, showmover=false, label=true, labelfontsize=6pt,
  boardfontsize=14pt}

% A move or line in UCI, and a position in FEN, both breakable at a slash.
\newcommand{\uci}[1]{\texttt{#1}}
\newcommand{\fen}[1]{{\footnotesize\path{#1}}}
\newcolumntype{L}{>{\raggedright\arraybackslash}X}

\setlength{\parskip}{0.4em}
\setlength{\parindent}{0pt}

% Written by projects/review-lessons/paper/generate.mjs; do not edit by hand.

\newcommand{\AppCommit}{816c07c243}
\newcommand{\Lessons}{74}
\newcommand{\Chapters}{493}
\newcommand{\WrongMoves}{114}
\newcommand{\WithMistakes}{41}
\newcommand{\TrainUntil}{2026-09-16}
\newcommand{\HeldFrom}{2026-09-17}
\newcommand{\HeldUntil}{2026-11-04}
\newcommand{\TrainEvery}{eight}
\newcommand{\Inspected}{2026-10-31, 2026-11-02, 2026-11-03}
\newcommand{\InspectedN}{three}
\newcommand{\JudgeTrain}{60}
\newcommand{\JudgeHeld}{120}
\newcommand{\Controls}{20}
\newcommand{\Draws}{10,000}
\newcommand{\SameBar}{16}
\newcommand{\DecoyBar}{14}
\newcommand{\CoverageSlack}{2}
\newcommand{\MistakePlies}{six}
\newcommand{\TrainGames}{59}
\newcommand{\TrainMoments}{334}
\newcommand{\TrainDiffer}{41}
\newcommand{\TrainBCoverage}{86.8}
\newcommand{\TrainCCoverage}{87.1}
\newcommand{\TrainBChapters}{18}
\newcommand{\TrainCChapters}{28}
\newcommand{\TrainBLessons}{17}
\newcommand{\TrainCLessons}{20}
\newcommand{\TrainBTop}{44.5}
\newcommand{\TrainCTop}{39.2}
\newcommand{\TrainViaMistake}{42}
\newcommand{\TrainViaFix}{171}
\newcommand{\TrainViaReply}{78}
\newcommand{\TrainViaNone}{43}
\newcommand{\TrainRead}{33}
\newcommand{\HeldGames}{223}
\newcommand{\HeldMoments}{1,225}
\newcommand{\HeldDiffer}{115}
\newcommand{\HeldBCoverage}{87.4}
\newcommand{\HeldCCoverage}{87.6}
\newcommand{\HeldBChapters}{24}
\newcommand{\HeldCChapters}{35}
\newcommand{\HeldBLessons}{21}
\newcommand{\HeldCLessons}{25}
\newcommand{\HeldBTop}{41.5}
\newcommand{\HeldCTop}{36.9}
\newcommand{\HeldViaMistake}{120}
\newcommand{\HeldViaFix}{623}
\newcommand{\HeldViaReply}{330}
\newcommand{\HeldViaNone}{152}
\newcommand{\HeldRead}{32}
\newcommand{\TopTitle}{Move the queen while taking the loose pawn}
\newcommand{\TopLesson}{Choose a defense that changes the threat}
\newcommand{\ReadWrong}{47}
\newcommand{\KindTable}{%
\begin{table}[ht]\centering\small
\begin{tabularx}{\textwidth}{lLrrrc}
\toprule
Kind & What the board code sees & Lessons & Train & Held out & C links \\
\midrule
\texttt{allows-mate} & the line mates within three of their moves & 16 & 15 & 40 & yes \\
\texttt{bad-capture} & his capture loses the exchange, and they take back on that square & 9 & 3 & 18 & yes \\
\texttt{unsafe-square} & their first move takes the moved piece and wins the exchange & 7 & 9 & 50 & no \\
\texttt{opens-line} & their first move captures along a line through the square his piece left & 4 & 5 & 7 & no \\
\texttt{deserted-guard} & the moved piece guarded a piece that now has no guard, and they take it within three plies & 5 & 10 & 25 & yes \\
\texttt{ignored-threat} & their first move takes a piece of his that was already loose & 4 & 22 & 94 & no \\
\texttt{allows-tactic} & a move of theirs in the line wins material with a tactic the board code names & 2 & 30 & 79 & yes \\
\texttt{loses-material} & the line loses two points or more, their capture within three plies & 0 & 16 & 82 & no \\
none & no kind fits, or the line loses no material & & 224 & 830 & \\
\bottomrule
\end{tabularx}
\caption{What his move got wrong, read from the move and the reply line the review's lesson shows, in the order the reader tries them at the capture it reads. \emph{Lessons}: the lessons' wrong moves read as that kind. \emph{Train}, \emph{Held out}: his moments. \emph{C links}: C names the chapter showing the same mistake for this kind; the others fall back to B.}
\label{tab:kinds}
\end{table}
}
\newcommand{\Kept}{kept}
\newcommand{\TrainJudged}{40}
\newcommand{\TrainWins}{30}
\newcommand{\TrainTies}{5}
\newcommand{\TrainLosses}{5}
\newcommand{\TrainNet}{0.63 (0.36 to 0.84)}
\newcommand{\HeldJudged}{113}
\newcommand{\HeldWins}{96}
\newcommand{\HeldTies}{7}
\newcommand{\HeldLosses}{10}
\newcommand{\HeldNet}{0.76 (0.65 to 0.86)}
\newcommand{\SameSame}{20}
\newcommand{\SameJudged}{20}
\newcommand{\DecoyWins}{14}
\newcommand{\DecoyJudged}{20}
\newcommand{\DecoySame}{4}
\newcommand{\DecoyLosses}{2}
\newcommand{\RuleTable}{%
\begin{table}[ht]\centering\small
\begin{tabular}{lrr}
\toprule
 & Train & Held out \\
\midrule
Games (moments) & 59 (334) & 223 (1,225) \\
Moments where B and C name different chapters & 41 & 115 \\
\quad judged (both name one) & 40 & 113 \\
\textbf{C wins / ties / losses} & \textbf{30 / 5 / 5} & \textbf{96 / 7 / 10} \\
\textbf{(wins $-$ losses) / judged} (bar: interval above 0) & 0.63 (0.36 to 0.84) & \textbf{0.76 (0.65 to 0.86)} \\
\textbf{Coverage}, B; C (bar: C at least B $-$ 2 points) & 86.8\%; 87.1\% & 87.4\%; \textbf{87.6\%} \\
\textbf{Most named chapter's share}, B; C (bar: C at most B) & 44.5\%; 39.2\% & 41.5\%; \textbf{36.9\%} \\
Distinct chapters (lessons), B; C & 18 (17); 28 (20) & 24 (21); 35 (25) \\
\bottomrule
\end{tabular}
\caption{The rule's measures. A moment is a win for C when the judge preferred C in both orders, or in one with the other the same; a loss likewise for B. Intervals are 95\% percentiles of 10,000 bootstrap draws over games. Coverage is the share of moments where the matcher names a chapter.}
\label{tab:rule}
\end{table}
}
\newcommand{\ByKindTable}{%
\begin{table}[ht]\centering\small
\begin{tabular}{lccc}
\toprule
His mistake & Train, six kinds & Train, four kinds & Held out, four kinds \\
\midrule
\texttt{allows-mate} & 10 / 0 / 2 & 13 / 0 / 2 & 34 / 2 / 4 \\
\texttt{deserted-guard} & 8 / 1 / 1 & 8 / 1 / 1 & 23 / 1 / 0 \\
\texttt{bad-capture} & 2 / 1 / 0 & 2 / 1 / 0 & 15 / 2 / 1 \\
\texttt{unsafe-square} & 1 / 3 / 4 & -- & -- \\
\texttt{opens-line} & 1 / 0 / 4 & -- & -- \\
\texttt{allows-tactic} & 7 / 3 / 2 & 7 / 3 / 2 & 24 / 2 / 5 \\
\midrule
All & 29 / 8 / 13 & 30 / 5 / 5 & 96 / 7 / 10 \\
\bottomrule
\end{tabular}
\caption{C's wins / ties / losses by what his move got wrong. The first column is the train round of the candidate before it was narrowed; the kinds C no longer links by fall back to B and leave the judged pairs.}
\label{tab:bykind}
\end{table}
}
\newcommand{\SixWins}{29}
\newcommand{\SixTies}{8}
\newcommand{\SixLosses}{13}
\newcommand{\SixJudged}{50}
\newcommand{\PathTable}{%
\begin{table}[ht]\centering\small
\begin{tabular}{lrrrrc}
\toprule
C's chapter, from & Held-out moments & C wins & same & C loses & Shown \\
\midrule
the same mistake (\texttt{mistake}) & 120 & 19 & 0 & 1 & yes \\
the engine's move (\texttt{fix}) & 623 & 15 & 3 & 2 & yes \\
their reply (\texttt{reply}) & 330 & 14 & 4 & 2 & yes \\
\bottomrule
\end{tabular}
\caption{Each way C names a chapter, against a chapter drawn at random from the other lessons: up to 20 held-out moments each, judged in both orders. A path is shown in the Game Review when C wins at least 14.}
\label{tab:paths}
\end{table}
}
\newcommand{\PathMistakeWins}{19}
\newcommand{\PathMistakeJudged}{20}
\newcommand{\PathFixWins}{15}
\newcommand{\PathFixJudged}{20}
\newcommand{\PathReplyWins}{14}
\newcommand{\PathReplyJudged}{20}
\newcommand{\EaSide}{Black}
\newcommand{\EaFen}{\fen{rn2k2r/pppb1pbp/4p1p1/1B4B1/8/2N2P2/PPP2P1P/2KR3R b kq - 3 12}}
\newcommand{\EaPlayed}{12\ldots{}\symbishop{}xb5}
\newcommand{\EaPlayedUci}{\uci{d7b5}}
\newcommand{\EaBest}{12\ldots{}h6}
\newcommand{\EaBestUci}{\uci{h7h6}}
\newcommand{\EaLine}{12\ldots{}\symbishop{}xb5 13.~\symrook{}d8\#}
\newcommand{\EaLineUci}{\uci{d7b5 d1d8}}
\newcommand{\EaKind}{\texttt{allows-mate} (\texttt{opera-mate})}
\newcommand{\EaBLesson}{Prove that the king has no escape}
\newcommand{\EaBChapter}{The checking queen is protected}
\newcommand{\EaCLesson}{Check what your move leaves behind}
\newcommand{\EaCChapter}{A bishop move that abandons d2}
\newcommand{\EaBoard}{%
\chessboard[setfen={rn2k2r/pppb1pbp/4p1p1/1B4B1/8/2N2P2/PPP2P1P/2KR3R b kq - 3 12}, boardfontsize=11pt, pgfstyle=straightmove, arrow=to, linewidth=0.07em, color=red!75!black, markmoves={d7-b5}, pgfstyle=straightmove, arrow=to, linewidth=0.07em, color=green!50!black, markmoves={h7-h6}]
}
\newcommand{\EbSide}{White}
\newcommand{\EbFen}{\fen{5r1k/p1p3b1/1p1pr2p/7n/4bp1B/2P2N1P/PP2BPRP/6RK w - - 0 28}}
\newcommand{\EbPlayed}{28.~b4}
\newcommand{\EbPlayedUci}{\uci{b2b4}}
\newcommand{\EbBest}{28.~\symrook{}f1}
\newcommand{\EbBestUci}{\uci{g1f1}}
\newcommand{\EbLine}{28.~b4 \symbishop{}b7 29.~\symbishop{}d1 \symbishop{}xc3}
\newcommand{\EbLineUci}{\uci{b2b4 e4b7 e2d1 g7c3}}
\newcommand{\EbKind}{\texttt{deserted-guard}}
\newcommand{\EbBLesson}{Improve with a tactic in reserve}
\newcommand{\EbBChapter}{The knight defends f7 and sets a trap}
\newcommand{\EbCLesson}{Check what your move leaves behind}
\newcommand{\EbCChapter}{One bishop move removes both guards}
\newcommand{\EbBoard}{%
\chessboard[setfen={5r1k/p1p3b1/1p1pr2p/7n/4bp1B/2P2N1P/PP2BPRP/6RK w - - 0 28}, boardfontsize=11pt, pgfstyle=straightmove, arrow=to, linewidth=0.07em, color=red!75!black, markmoves={b2-b4}, pgfstyle=straightmove, arrow=to, linewidth=0.07em, color=green!50!black, markmoves={g1-f1}]
}
\newcommand{\EcSide}{White}
\newcommand{\EcFen}{\fen{r3k1n1/2Bb4/p3p1p1/1p1p1pq1/PQ1P4/2PBr3/1P4P1/R3R1K1 w q - 0 21}}
\newcommand{\EcPlayed}{21.~\symrook{}xe3}
\newcommand{\EcPlayedUci}{\uci{e1e3}}
\newcommand{\EcBest}{21.~\symbishop{}f1}
\newcommand{\EcBestUci}{\uci{d3f1}}
\newcommand{\EcLine}{21.~\symrook{}xe3 \symqueen{}xe3+ 22.~\symking{}h1 \symqueen{}xd3}
\newcommand{\EcLineUci}{\uci{e1e3 g5e3 g1h1 e3d3}}
\newcommand{\EcKind}{\texttt{allows-tactic} (\texttt{fork})}
\newcommand{\EcBLesson}{Choose a defense that changes the threat}
\newcommand{\EcBChapter}{Move the queen while taking the loose pawn}
\newcommand{\EcCLesson}{Recognize tactics during an ordinary game}
\newcommand{\EcCChapter}{A knight attacks king, queen, and rook}
\newcommand{\EcBoard}{%
\chessboard[setfen={r3k1n1/2Bb4/p3p1p1/1p1p1pq1/PQ1P4/2PBr3/1P4P1/R3R1K1 w q - 0 21}, boardfontsize=11pt, pgfstyle=straightmove, arrow=to, linewidth=0.07em, color=red!75!black, markmoves={e1-e3}, pgfstyle=straightmove, arrow=to, linewidth=0.07em, color=green!50!black, markmoves={d3-f1}]
}


\title{Which lesson chapter should a Game Review moment link to?\\[0.3em]
\large The chapter that shows the same mistake}
\author{Chess Coach research lab}
\date{}

\begin{document}
\maketitle

\begin{abstract}
The app's Game Review explains each key moment of a game, but it named no lesson, although the
app has \Lessons{} study lessons. The daily review does name one: it reads what the engine's
better move is for, with the board code, and offers the chapter whose recommended moves are read
the same way (B). That matcher concentrates: one chapter takes \HeldBTop\% of the held-out moments
it names. The lessons also play \WrongMoves{} wrong moves on purpose. The candidate (C) reads what
his move got wrong on the board, from the move and the reply line the review's lesson uses, and
for four kinds of mistake links the chapter whose wrong move is read the same way; otherwise it is
B. On \HeldGames{} held-out games, a blind judge, in both orders, preferred C at \HeldWins, tied
\HeldTies{} and preferred B at \HeldLosses{} of the \HeldJudged{} moments where both name a
chapter and the chapters differ: (wins $-$ losses) / judged $=$ \HeldNet. The controls held, coverage stayed,
and the most named chapter's share fell to \HeldCTop\%. By the rule fixed before the run, C is
\Kept.
\end{abstract}

\section{The question}

At a key moment where he went wrong, the review says what his move allowed and what the
engine's move does instead. A link to the lesson that teaches the idea would turn that into
study. The daily review already links one, by what the engine's move is for, and on his games
one chapter, ``\TopTitle'' from ``\TopLesson'', is its pick at \HeldBTop\% of the held-out
moments it names a chapter for, whatever his move got wrong. The lessons hold another source: the wrong moves they play on purpose,
each in a chapter about that mistake. Can the board code read his mistake the way those chapters
show it, and does a link to the chapter showing the same mistake fit the moment better?

\section{What the app does}

\paragraph{What his move got wrong.} The Knowledge service's \texttt{mistakeOf} walks his move
and the reply line the review's lesson shows after it (Table~\ref{tab:kinds}). A mate comes
first: the line mates within three of their moves, with the mate's pattern when the board code
names it. Otherwise the line must lose material. The reader takes their first capture in the
first \MistakePlies{} plies, passing over their taking of the moved piece when that exchange does
not lose for him, and names the first kind that fits that capture, or none.

\KindTable

\paragraph{The lessons' wrong moves.} The scripts mark \texttt{counterexample} each wrong move
the lessons play on purpose, and the study build reads them the same way, with the scene's moves
after each as the line: \ReadWrong{} of the \WrongMoves{} are read as a kind, in \WithMistakes{} of the \Chapters{}
chapters. The others lose nothing in the scene's line, or lose it in a way no kind names.

\paragraph{The match.} For his mistake, a chapter is offered when one of its wrong moves is read
as the same kind, with the same idea (the mate's pattern, the tactic) or with no idea named, and
its lesson is tagged with one of the mistake's concepts; a wrong move set in the opening counts
only for a mistake in the opening. The same idea counts most, then the same phase, the lesson's topic, the
same piece moved and the same piece lost. When no chapter shows the mistake and it has an idea,
the chapters whose own recommended moves play that idea, in lessons that list it among their
first two tags, are offered instead.

\paragraph{The two matchers.} B is the daily review's: the first chapter offered for what the
engine's move is for, else for what their reply is for. C is the chapter matched to his mistake
when it is one of the four kinds the last column of Table~\ref{tab:kinds} marks, else B. The code
is the app's at commit \texttt{\AppCommit}.

\section{How it was measured}

\paragraph{The moments.} His rapid games in the app's archive, each reviewed the way the phone's
Game Review reviews it. A moment is a key moment of the review where he moved, the move is graded
an inaccuracy, mistake, miss or blunder, and the review writes a lesson for it.

\paragraph{The split.} Train: one game in every \TrainEvery{} up to \TrainUntil{} (\TrainGames{}
games, \TrainMoments{} moments), on which the change was designed. Held out: every game from
\HeldFrom{} to \HeldUntil{} (\HeldGames{} games, \HeldMoments{} moments) except those of the
\InspectedN{} days read while it was designed (\Inspected).

\paragraph{The judge.} On the moments where B and C name different chapters (at most
\JudgeTrain{} on train and \JudgeHeld{} held out), a model judge sees the moment as the review
describes it (the position's FEN, his move and its grade, the engine's move, the reply line, the
lesson's headline and its ``difference'' lines) and the two chapters, each as its lesson's title,
its chapter's title and its narration, labelled A and B. It answers which chapter teaches what
the moment needed, or that they are the same. Every pair is judged in both orders.

\paragraph{The rule, fixed before the run.} C is kept when, held out, (wins $-$ losses) / judged
has its 95\% interval above zero, from \Draws{} bootstrap draws over games; C's coverage is at
least B's minus \CoverageSlack{} points; C's most named chapter takes no larger share than B's;
on train, C's wins are not fewer than its losses; and the controls hold: \Controls{} held-out
moments with the same chapter on both sides are read as the same at least \SameBar{} times, and
\Controls{} with C's chapter against a chapter drawn from another lesson are won by C at least
\DecoyBar{} times.

\paragraph{Changed after the train round.} C first linked by the mistake for six kinds. On train
(C \SixWins{} wins, \SixTies{} ties, \SixLosses{} losses of \SixJudged), the judge favoured B on
the \texttt{unsafe-square} and \texttt{opens-line} moments and C on every other kind
(Table~\ref{tab:bykind}, first column), so those two now fall back to B. The change was made
before any held-out verdict, and the rule was applied unchanged to the narrowed C.

\section{The verdict}

\RuleTable

\ByKindTable

On held-out moments C wins \HeldWins{} and loses \HeldLosses{} of \HeldJudged{}
(Table~\ref{tab:rule}), the controls hold (the same chapter read as the same \SameSame{} of
\SameJudged{} times; C beat the decoy \DecoyWins{} of \DecoyJudged{} times), C names a chapter
about as often as B, and it spreads the links over more chapters: \HeldCChapters{} against
\HeldBChapters. C is \Kept.

\paragraph{Where the link comes from.} Held out, C names the chapter showing the same mistake at
\HeldViaMistake{} moments, the chapter for what the engine's move is for at \HeldViaFix, the
chapter for what their reply is for at \HeldViaReply, and none at \HeldViaNone. Beside the rule,
each path was judged against a decoy chapter from another lesson (Table~\ref{tab:paths}); the
Game Review shows a link only from a path that reaches the bar.

\PathTable

\clearpage
\section{Three moments}

\begin{figure}[!ht]
\centering
\begin{subfigure}[t]{0.31\textwidth}
\centering
\EaBoard\\[0.4em]
{\footnotesize FEN \EaFen}
\caption{\EaSide{} to move. He played \EaPlayed{} (\EaPlayedUci, red); the engine preferred
\EaBest{} (\EaBestUci, green). The line: \EaLine{} (\EaLineUci). Read as \EaKind. B: ``\EaBChapter''
from ``\EaBLesson''. C: ``\EaCChapter'' from ``\EaCLesson''.}
\end{subfigure}\hfill
\begin{subfigure}[t]{0.31\textwidth}
\centering
\EbBoard\\[0.4em]
{\footnotesize FEN \EbFen}
\caption{\EbSide{} to move. He played \EbPlayed{} (\EbPlayedUci); the engine preferred \EbBest{}
(\EbBestUci). The line: \EbLine{} (\EbLineUci). Read as \EbKind. B: ``\EbBChapter'' from
``\EbBLesson''. C: ``\EbCChapter'' from ``\EbCLesson''.}
\end{subfigure}\hfill
\begin{subfigure}[t]{0.31\textwidth}
\centering
\EcBoard\\[0.4em]
{\footnotesize FEN \EcFen}
\caption{\EcSide{} to move. He played \EcPlayed{} (\EcPlayedUci); the engine preferred \EcBest{}
(\EcBestUci). The line: \EcLine{} (\EcLineUci). Read as \EcKind. B: ``\EcBChapter'' from
``\EcBLesson''. C: ``\EcCChapter'' from ``\EcCLesson''.}
\end{subfigure}
\caption{Held-out moments where the judge preferred C in both orders, one for each of three kinds.}
\label{fig:moments}
\end{figure}

\section{Limits}

The judge is a language model reading the review's own words and the chapters' narration; it is
told to take the review's facts as true and not to analyse the position, so it judges fit, not
chess. Its fairness is checked only by the two controls, and the decoy control passed with C
winning \DecoyWins{} of \DecoyJudged, where \DecoyBar{} were needed. By path
(Table~\ref{tab:paths}), the chapter showing the same mistake beat its decoy \PathMistakeWins{} of
\PathMistakeJudged{} times, the chapters B's two readings name \PathFixWins{} and
\PathReplyWins{} times: the links that come from B are the weaker part of C. The moments are one player's, at his
level. The mistake reader names a kind from one reply line, the one the review shows; a move can
be wrong in more ways than that line shows. A kind with few wrong moves in the lessons links few
chapters, so C's links stay concentrated where the lessons show only one example.

\section{What the app took}

The Game Review's lesson panel, at his key moments, ends with the chapter C names, from the paths
Table~\ref{tab:paths} marks as shown, as a link to watch or read the lesson with a line on why it
fits. The study build publishes the chapters with what each one's moves are for and get wrong, so
the review links one without loading every lesson.

\end{document}
